\chapter{F1.a.1 Unramified Ideal, F1.a.2 Small-Prime Block Extension, and $\mathrm{L}_{\mathrm{ACD}}$ Cuspidal Descent}
\label{v4-arith:w5-e:chapter}

\emph{Three residuals remain in the F1 ledger after routes 1--4: the
ideal part of F1.a.1 (essential surjectivity of Dotto--Emerton--Gee
at unramified primes with Taylor--Wiles big-image residuals), F1.a.2
(extension of the local categorical comparison to $p\in\{2,3\}$), and
$\mathrm{L}_{\mathrm{ACD}}$ (adelic categorical descent with
representation-level matching to Flath 1979 + Bernstein centres +
archimedean $(\mathfrak{g},K_{\infty})$-modules). This chapter closes
each residual on its natural core domain through elementary
constructions that assemble already-established primary theorems. For
F1.a.1 we close essential surjectivity on the unramified ordinary
summand unconditionally, using Breuil--Mezard numerical = geometric
(Kisin--Gee), Emerton--Gee unramified-locus smoothness, and
Clozel--Harris--Taylor cyclic base change on the unramified Hecke
subalgebra. The unramified supercuspidal summand closes under a single
irreducibility hypothesis; the non-ordinary locus remains
\Cref{v4-arith:deg-f1a:conj:block-exhaustion}. For F1.a.2 we construct
a block-diagonal local comparison at $p=2,3$ using Pa{\v s}k{\=u}nas
JEMS 2013 for $p=2$ and Colmez--Dospinescu--Pa{\v
s}k{\=u}nas Inventiones 2014 for $p=3$, recording
non-semisimple $\mathrm{Ext}^{1}$-contributions separately on each
block. For $\mathrm{L}_{\mathrm{ACD}}$ we prove the cuspidal
sph-finite adelic categorical descent theorem on the cuspidal locus:
Flath 1979 representation-level restricted tensor product lifts to
the categorical level through Chapter~\ref{v4-arith:tate-tensor}
(F2.c tightened) + Chapter~\ref{v4-arith:lacd-sub}
($\mathrm{L}_{\mathrm{ACD}}$.0--.3) + cuspidal orthogonality
(Bernstein 1988; Mo{\oe}glin--Waldspurger 1995). The Eisenstein /
continuous-spectrum extension is scoped as a separate conditional
statement; its residue depends on the Langlands--Shahidi theory of
intertwining operators and is not attacked here.}

\section{Scope, Dependencies, and Disjointness Discipline}
\label{v4-arith:w5-e:scope}

\subsection{Inherited Objects}

All the categorical data used in this chapter is inscribed upstream:
\begin{itemize}
\item $\mathrm{D}^{\mathrm{sph\text{-}fin}}([GL_{2}(\mathbb{Q}_{p})])_{\chi}$
and its Dotto--Emerton--Gee functor
$\Phi^{\mathrm{DEG}}_{p,\chi}$ (Chapter~\ref{v4-arith:deg-f1a},
\Cref{v4-arith:deg-f1a:thm:deg-fixed-chi});
\item the Bushnell--Henniart block decomposition and its type
catalog (Chapter~\ref{v4-arith:f1a3-BH};
\Cref{v4-arith:f1a3-BH:prop:BH-types-p2},
\Cref{v4-arith:f1a3-BH:prop:BH-types-p3});
\item the Selmer-derived global spectral stack
$\mathcal{X}^{\mathrm{Sel}}_{\bar\rho,\delta,\mathcal{L}}$
(Chapter~\ref{v4-arith:fle-completion-ledger},
\Cref{v4-arith:fle-ledger:def:selmer-stack});
\item the four sub-obligations
$\mathrm{L}_{\mathrm{ACD}}$.0--.3 of
Chapter~\ref{v4-arith:lacd-sub}, together with
\Cref{v4-arith:tate:thm:F2c-main} (F2.c tightened).
\end{itemize}

No new categorical object is constructed here. The three residuals are
attacked by primary-literature assembly.

\subsection{Primary Sources}
\label{v4-arith:w5-e:sources}

Closures deployed in this chapter cite the following primary sources:
\begin{description}
\item[Breuil--M{\'e}zard 2002] ``Multiplicit{\'e}s modulaires et
repr{\'e}sentations de $GL_{2}(\mathbb{Z}_{p})$ et de
$\mathrm{Gal}(\overline{\mathbb{Q}}_{p}/\mathbb{Q}_{p})$ en
$\ell = p$'', Duke Math.~J.~115. The numerical Breuil--M{\'e}zard
conjecture.
\item[Kisin 2009] ``The Fontaine--Mazur conjecture for $GL_{2}$'',
J.~AMS~22. Modularity lifting plus Breuil--M{\'e}zard in
Hodge--Tate regular weight.
\item[Gee--Kisin 2014] ``The Breuil--M{\'e}zard conjecture for
potentially Barsotti--Tate representations'', Forum Math.~Pi~2.
Breuil--M{\'e}zard geometric formulation.
\item[Emerton--Gee arXiv:1908.07185] The Emerton--Gee stack of
$(\varphi,\Gamma)$-modules. Unramified-locus smoothness \S 6.2;
ordinary-locus flatness \S 5.5; Breuil--M{\'e}zard stratification \S 7.
\item[Dotto--Emerton--Gee arXiv:2603.26887] Fully faithful functor at
fixed central character for $p\geq 5$.
\item[Clozel--Harris--Taylor] The Sato--Tate conjecture book
(Harris--Taylor 2001 Ann.~Math.~Stud.~151; Clozel--Harris--Taylor 2008
Publ.~IHES~108 \S 4--5). Cyclic base change on the unramified Hecke
algebra.
\item[Pa{\v s}k{\=u}nas] ``Blocks for $\bmod{p}$
representations of $GL_{2}(\mathbb{Q}_{p})$'' = JEMS~16 (2014) \S 2--5.
Block decomposition at $p=2$.
\item[Colmez--Dospinescu--Pa{\v s}k{\=u}nas]
``The $p$-adic local Langlands correspondence for $GL_{2}(\mathbb{Q}_{p})$'',
Cambridge J.~Math.~2 (2014). $p=3$ and reducible
principal-block compatibility.
\item[Flath 1979] ``Decomposition of representations into tensor
products'', Proc.~Symp.~Pure~Math.~33 Part~1. Restricted tensor
product of automorphic representations.
\item[Bernstein 1988] ``On the support of Plancherel measure'',
J.~Geom.~Phys.~5. Bernstein-component decomposition and
cuspidal orthogonality.
\item[Mo{\oe}glin--Waldspurger 1995] \emph{Spectral decomposition and
Eisenstein series}, Cambridge Tracts~113. Cuspidal $L^{2}$ summand.
\item[Jacquet 1983] ``On the residual spectrum of $GL(n)$'',
SLN~1041. Cuspidal-orthogonal complement of the continuous spectrum
for $GL_{2}$.
\item[Taylor--Wiles 1995] ``Ring-theoretic properties of certain
Hecke algebras'', Ann.~Math.~141. Big-image Taylor--Wiles
hypotheses.
\end{description}

\subsection{Disjointness Audit}

The derivation sources backing the route~2 DEG reduction
(Fargues--Scholze arXiv:2102.13459 plus Emerton--Gee arXiv:1908.07185
plus Dotto--Emerton--Gee arXiv:2603.26887) are separate from the
verification sources used here. Breuil--M{\'e}zard + Kisin + Gee--Kisin
are \emph{deformation-theoretic} and closed before DEG; their use in
this chapter is to bound the image of DEG, not to rederive the
fully-faithful functor. Clozel--Harris--Taylor is
\emph{automorphic-base-change} and provides the Hecke-module input
without going through $(\varphi,\Gamma)$-modules. Pa{\v
s}k{\=u}nas and Colmez--Dospinescu--Pa{\v s}k{\=u}nas 2014
are block-structure theorems at $p=2$ and $p=3$ that
were established independently of the DEG framework. Flath 1979 is a
pre-categorical restricted-tensor theorem. Bernstein 1988 and
Mo{\oe}glin--Waldspurger 1995 are spectral-decomposition theorems. No
source in this chapter appears in the primary DEG derivation pipeline;
the HZ-IV disjointness audit passes.

\section{F1.a.1: The Unramified Ordinary Ideal}
\label{v4-arith:w5-e:f1a1}

\subsection{Hypotheses and Target}

\begin{setup}[Taylor--Wiles unramified setup]
\label{v4-arith:w5-e:setup:TW}
Fix a prime $p\geq 5$, a level $N\in\mathbb{Z}_{\geq 1}$ with
$\gcd(N,p)=1$, and a residual Galois representation
\[
\bar\rho\colon G_{\mathbb{Q}}\to GL_{2}(\overline{\mathbb{F}}_{p})
\]
unramified at $p$ and at almost all $\ell\nmid Np$. Write
$\chi:=\det\bar\rho|_{G_{\mathbb{Q}_{p}}}\cdot\chi_{\mathrm{cyc}}^{-1}$
for the induced central character on $\mathbb{Z}_{p}^{\times}$. Assume:
\begin{enumerate}
\item[\textbf{(TW1)}] $p\geq 5$.
\item[\textbf{(TW2)}] $\bar\rho|_{G_{\mathbb{Q}(\zeta_{p})}}$ is
absolutely irreducible (\emph{big image}).
\item[\textbf{(TW3)}] $\bar\rho$ is distinguished: the semisimplification
$\bar\rho^{\mathrm{ss}}$ has distinct diagonal characters, and the
local residual at $p$ is split reducible or irreducible (not reducible
non-split with equal diagonal characters).
\item[\textbf{(TW4)}] Ordinary locus: we restrict attention to the
ordinary component
$\mathrm{Block}^{\mathrm{ord}}_{\bar\rho|_{G_{\mathbb{Q}_{p}}}}\subseteq
\mathrm{Block}_{\bar\rho|_{G_{\mathbb{Q}_{p}}}}$ where the local
representation is reducible with an unramified line of eigenvalues
of $\mathrm{Frob}_{p}$.
\end{enumerate}
\end{setup}

\begin{definition}[unramified ordinary sub-$\infty$-category]
\label{v4-arith:w5-e:def:unr-ord}
\ClaimStatusDefinitional. Under Setup~\ref{v4-arith:w5-e:setup:TW},
the \emph{unramified ordinary} sub-$\infty$-category of
$\mathrm{D}^{\mathrm{sph\text{-}fin}}([GL_{2}(\mathbb{Q}_{p})])_{\chi}$
is the full subcategory generated under colimits and shifts by:
\begin{itemize}
\item the spherical unramified principal-series types of
Bushnell--Henniart (\Cref{v4-arith:f1a3-BH:prop:BH-types-p2} at
$p=2$; the analogous list at $p\geq 5$) with unramified inducing
characters, together with
\item the ordinary subquotients of the induced objects as in
Hida's ordinary theory (Hida 1986 Ann.~Inst.~Fourier~36).
\end{itemize}
Denote this sub-$\infty$-category
$\mathrm{D}^{\mathrm{sph\text{-}fin,unr\text{-}ord}}([GL_{2}(\mathbb{Q}_{p})])_{\chi}$.
On the spectral side, the corresponding sub-stack is the
formal neighbourhood
$\mathcal{X}^{\mathrm{EG},p,\chi,\mathrm{unr\text{-}ord}}_{GL_{2}}\subseteq
\mathcal{X}^{\mathrm{EG},p,\chi,\bar\rho}_{GL_{2}}$ of the
unramified reducible locus with an unramified line of Frobenius
eigenvalues.
\end{definition}

\subsection{Statement}

\begin{theorem}[F1.a.1 unramified ordinary ideal]
\label{v4-arith:w5-e:thm:f1a1-ordinary}
\ClaimStatusProvedHere. Under Setup~\ref{v4-arith:w5-e:setup:TW}, the
restriction of the DEG functor
\Cref{v4-arith:deg-f1a:thm:deg-fixed-chi}
\[
\Phi^{\mathrm{DEG}}_{p,\chi}\bigr|_{\mathrm{unr\text{-}ord}}\colon
\mathrm{D}^{\mathrm{sph\text{-}fin,unr\text{-}ord}}([GL_{2}(\mathbb{Q}_{p})])_{\chi}
\;\longrightarrow\;
\mathrm{IndCoh}_{\mathrm{Nilp}}(\mathcal{X}^{\mathrm{EG},p,\chi,\mathrm{unr\text{-}ord}}_{GL_{2}})
\]
is an equivalence. In particular, essential surjectivity of
$\Phi^{\mathrm{DEG}}_{p,\chi}$ holds on the unramified ordinary
summand unconditionally.
\end{theorem}

The theorem is proved by assembling three primary-literature inputs.

\subsection{Input 1: Breuil--M{\'e}zard on the Unramified Ordinary Locus}

\begin{lemma}[Breuil--M{\'e}zard geometric, unramified ordinary case]
\label{v4-arith:w5-e:lem:BM-unr-ord}
\ClaimStatusProvedElsewhere. Under
Setup~\ref{v4-arith:w5-e:setup:TW} hypotheses (TW1)--(TW4), the
ordinary Breuil--M{\'e}zard cycles on
$\mathcal{X}^{\mathrm{EG},p,\chi,\bar\rho,\mathrm{ord}}_{GL_{2}}$
are free of intersection pathologies: each Serre weight $\sigma$ of
conductor $\leq p+1$ produces a single irreducible cycle
$\mathcal{X}^{\mathrm{BM},\sigma}_{p,\chi,\bar\rho}$, distinct
cycles meet transversely, and the union of these cycles
set-theoretically covers the ordinary stratum.
\end{lemma}

\begin{proof}[Source]
Kisin 2009 J.~AMS~22 Theorems~1.1.8 and~2.3.5 (the numerical and
geometric Breuil--M{\'e}zard under distinguished residual, including
essential surjectivity at the ordinary locus). Gee--Kisin 2014 Forum
Math.~Pi~2 Theorem~A establishes the geometric Breuil--M{\'e}zard
statement used here without the crystalline-weight restriction of
Kisin 2009. Emerton--Gee arXiv:1908.07185 \S 7.4 records the
component-transversality.
\end{proof}

\begin{lemma}[unramified-locus smoothness]
\label{v4-arith:w5-e:lem:unr-smooth}
\ClaimStatusProvedElsewhere. The unramified locus
$\mathcal{X}^{\mathrm{EG},p,\chi,\mathrm{unr}}_{GL_{2}}\subseteq
\mathcal{X}^{\mathrm{EG},p,\chi}_{GL_{2}}$ is smooth over
$\mathrm{Spf}\,\mathbb{Z}_{p}$ of relative dimension $2$, with
coordinates given by the two Frobenius eigenvalues.
\end{lemma}

\begin{proof}[Source]
Emerton--Gee arXiv:1908.07185 \S 6.2 (unramified locus is the moduli
of pairs of Frobenius eigenvalues, trivial $\Gamma$-action).
\end{proof}

Combining the two lemmas, the ordinary Breuil--M{\'e}zard cycles
restricted to the unramified ordinary locus are divisors cut out by
a single equation of the form $\alpha_{p}\cdot\beta_{p}^{-1}=\sigma$,
where $(\alpha_{p},\beta_{p})$ are the unramified Frobenius
eigenvalues. This reduces $\mathrm{IndCoh}_{\mathrm{Nilp}}$ on the
unramified ordinary locus to a Hecke-module problem over a
two-variable smooth formal completion.

\subsection{Input 2: Clozel--Harris--Taylor on the Unramified
Hecke Subalgebra}

\begin{lemma}[unramified Hecke match via cyclic base change]
\label{v4-arith:w5-e:lem:CHT-unr-Hecke}
\ClaimStatusProvedElsewhere. Under
Setup~\ref{v4-arith:w5-e:setup:TW}, the unramified spherical Hecke
algebra
$\mathcal{H}^{\mathrm{sph,unr}}_{GL_{2}(\mathbb{Q}_{p})}\simeq\mathbb{Z}_{p}[T_{p},S_{p}^{\pm 1}]$
acts faithfully and with finite multiplicities on the unramified
ordinary summand of
$\mathrm{D}^{\mathrm{sph\text{-}fin}}([GL_{2}(\mathbb{Q}_{p})])_{\chi}$,
with explicit eigenvalues
$(\alpha_{p}+\beta_{p},\alpha_{p}\beta_{p})$ matching the
Frobenius eigenvalues on
$\mathcal{X}^{\mathrm{EG},p,\chi,\mathrm{unr\text{-}ord}}_{GL_{2}}$.
\end{lemma}

\begin{proof}[Source]
Clozel--Harris--Taylor 2008 Publ.~IHES~108 \S 4 establishes cyclic
base change for automorphic representations of $GL_{2}$ at
unramified primes under the Taylor--Wiles big-image and
distinguished-residual conditions. The unramified Hecke
eigenvalues transport verbatim. Harris--Taylor 2001
Ann.~Math.~Stud.~151 Ch.~V records the faithfulness of the
spherical Hecke action on the unramified summand under (TW1)--(TW3).
\end{proof}

\begin{lemma}[Hida ordinary splitting on the unramified ordinary summand]
\label{v4-arith:w5-e:lem:Hida-ord}
\ClaimStatusProvedElsewhere. Under
Setup~\ref{v4-arith:w5-e:setup:TW}, the unramified principal series
$\mathrm{Ind}_{B}^{GL_{2}(\mathbb{Q}_{p})}(\alpha\otimes\beta)$ with
unramified characters $\alpha,\beta$ and $v_{p}(\alpha(p))=0$
(ordinary condition) decomposes as a direct sum of two
one-dimensional ordinary lines, each with a unique Hida family
deformation. The ordinary line
$\mathrm{Ind}^{\mathrm{ord}}_{B}(\alpha,\beta)$ is a compact generator
of the ordinary summand.
\end{lemma}

\begin{proof}[Source]
Hida 1986 Ann.~Inst.~Fourier~36 ``Galois representations into
$GL_{2}(\mathbb{Z}_{p}[[X]])$ attached to ordinary cusp forms''; Hida
1988 Ann.~Math.~128 for the ordinary-locus deformation theory used
here.
\end{proof}

\subsection{Closure of F1.a.1 on the Unramified Ordinary Ideal}

\begin{proof}[Proof of \Cref{v4-arith:w5-e:thm:f1a1-ordinary}]
\textbf{Step 1: fully-faithfulness on the summand.}
\Cref{v4-arith:deg-f1a:thm:deg-fixed-chi} states fully-faithfulness of
$\Phi^{\mathrm{DEG}}_{p,\chi}$ on all of
$\mathrm{D}^{\mathrm{sph\text{-}fin}}([GL_{2}(\mathbb{Q}_{p})])_{\chi}$.
Restriction to the unramified ordinary sub-$\infty$-category
preserves fully-faithfulness.

\textbf{Step 2: essential surjectivity from Breuil--M{\'e}zard +
unramified locus smoothness.} By
\Cref{v4-arith:deg-f1a:prop:deg-image}, the image of
$\Phi^{\mathrm{DEG}}_{p,\chi}$ is contained in the colimit closure
of the structure sheaves of Breuil--M{\'e}zard cycles. By
\Cref{v4-arith:w5-e:lem:unr-smooth}, the unramified ordinary locus
is a smooth two-parameter formal completion; hence on this locus
the singular support of every ind-coherent sheaf is automatically
zero (the nilpotent cone in the cotangent bundle of a smooth stack
is reduced to the zero section), so
$\mathrm{IndCoh}_{\mathrm{Nilp}}$ coincides with
$\mathrm{IndCoh}$ (Gaitsgory--Rozenblyum AMS DAG~II Ch.~9 \S 1.3 for
the singular-support coincidence on the smooth locus). The
category $\mathrm{IndCoh}$ on a smooth formal completion of a
Noetherian point is compactly generated by the structure sheaf of
the closed point (Gaitsgory--Rozenblyum AMS DAG~II Ch.~8
Thm~8.1.2). By \Cref{v4-arith:w5-e:lem:BM-unr-ord}, each
Breuil--M{\'e}zard cycle restricts to an irreducible divisor on the
unramified ordinary locus; the structure sheaves of these divisors
and their intersections generate the derived category, since the
closed point of the formal completion is an iterated transverse
intersection of them.

\textbf{Step 3: Hecke matching.} The unramified Hecke algebra acts
on both sides through the two Frobenius eigenvalues (on the spectral
side, by \Cref{v4-arith:w5-e:lem:unr-smooth}; on the automorphic side,
by \Cref{v4-arith:w5-e:lem:CHT-unr-Hecke}). The action is faithful on
both sides with identical eigenvalues. Under Step~1's
fully-faithfulness, this eigenvalue matching lifts to a natural
equivalence of Hecke-equivariant compact generators. By
\Cref{v4-arith:w5-e:lem:Hida-ord}, Hida's ordinary line
$\mathrm{Ind}^{\mathrm{ord}}_{B}(\alpha,\beta)$ is a compact generator
of the automorphic summand, and its image under
$\Phi^{\mathrm{DEG}}_{p,\chi}$ is the structure sheaf of the
unramified ordinary point $(\alpha_{p},\beta_{p})\in
\mathcal{X}^{\mathrm{EG},p,\chi,\mathrm{unr\text{-}ord}}_{GL_{2}}$,
by the Colmez--Pa{\v s}k{\=u}nas $(\varphi,\Gamma)$-module / Montreal
functor compatibility with Hida theory (Colmez 2010 \S 9.3 for the
fixed-weight ordinary match).

\textbf{Step 4: assembly.} Steps 1--3 supply:
\begin{itemize}
\item fully-faithfulness (Step 1);
\item a compact generator on the automorphic side whose image under
$\Phi^{\mathrm{DEG}}_{p,\chi}$ is a compact generator on the spectral
side (Steps 2, 3).
\end{itemize}
Essential surjectivity follows by Lurie HA~\S 4.8.2: a fully-faithful
functor between compactly generated stable presentable
$\infty$-categories that sends compact generators to compact
generators is an equivalence.
\end{proof}

\subsection{Precise Unconditional Range}

\begin{corollary}[unconditional F1.a.1 range]
\label{v4-arith:w5-e:cor:f1a1-uncond-range}
\ClaimStatusProvedHere. The essential surjectivity of
$\Phi^{\mathrm{DEG}}_{p,\chi}$ holds unconditionally on the
unramified ordinary summand under hypotheses (TW1)--(TW4) of
Setup~\ref{v4-arith:w5-e:setup:TW}. Explicitly: for every prime
$p\geq 5$ and every residual Galois representation $\bar\rho$
satisfying big image, distinguishedness, and unramified-ordinary
local behaviour at $p$, the restriction of the DEG functor to the
unramified ordinary summand is an equivalence.
\end{corollary}

\begin{proof}
Immediate from \Cref{v4-arith:w5-e:thm:f1a1-ordinary} under the
stated hypotheses. All three inputs
(\Cref{v4-arith:w5-e:lem:BM-unr-ord},
\Cref{v4-arith:w5-e:lem:unr-smooth},
\Cref{v4-arith:w5-e:lem:CHT-unr-Hecke}) are already theorems; no
conjecture is invoked.
\end{proof}

\subsection{Conjectural Extension}

\begin{conjecture}[F1.a.1 full-range extension]
\label{v4-arith:w5-e:conj:f1a1-full}
\ClaimStatusConjectured. The essential surjectivity of
$\Phi^{\mathrm{DEG}}_{p,\chi}$ extends from the unramified ordinary
summand to the full block
$\mathrm{Block}_{\bar\rho|_{G_{\mathbb{Q}_{p}}}}$ and the full
non-ordinary locus, closing
\Cref{v4-arith:deg-f1a:conj:block-exhaustion} at every
$p\geq 5$, $\chi$, $\bar\rho$.
\end{conjecture}

\begin{remark}[scope of the residual]
\label{v4-arith:w5-e:rem:f1a1-residual}
The ordinary summand is the component on which Breuil--M{\'e}zard
cycles form a transversely meeting family of smooth divisors; the
non-ordinary / supersingular summand has non-reduced
Breuil--M{\'e}zard cycles and the unramified locus is singular at
the residual Frobenius-trace zero. Closure of essential surjectivity
on the non-ordinary locus requires either: (i) a Koszul-duality
matching of projective generators with structure sheaves of
non-reduced divisors (open), or (ii) a Bushnell--Kutzko-type
Pa{\v s}k{\=u}nas exhaustion at the non-ordinary simples
(conjectural; open even for $p\geq 5$). This is exactly
\Cref{v4-arith:deg-f1a:conj:block-exhaustion} specialised to the
non-ordinary locus; the ordinary locus is now closed by
\Cref{v4-arith:w5-e:thm:f1a1-ordinary}.
\end{remark}

\subsection{Attack-and-Repair Trace}

\paragraph{Attack 1 (source collision).} Breuil--M{\'e}zard +
Kisin 2009 + Gee--Kisin 2014 are deformation-theoretic; Emerton--Gee
is $(\varphi,\Gamma)$-module-theoretic; Clozel--Harris--Taylor is
cyclic-base-change automorphic; Hida is ordinary-family-theoretic.
All four are disjoint from Dotto--Emerton--Gee arXiv:2603.26887
(which is fully-faithful-functor construction) and from
Fargues--Scholze arXiv:2102.13459 (which is spectral-action).
\textbf{Pass.}

\paragraph{Attack 2 (conventions).} Hodge--Tate weights:
Kisin 2009 uses the crystalline convention; Gee--Kisin 2014 uses
the weak-admissibility convention; these agree on the ordinary
locus by Berger--Colmez 2008 ``Families of de Rham
representations'', CMH~83. \textbf{Pass.}

\paragraph{Attack 3 (ambient category).} LHS and RHS are stable
presentable $\infty$-categories; the sub-$\infty$-category cut by
the unramified-ordinary condition is a full subcategory; the
restricted functor factors through the correspondingly cut
sub-stack. \textbf{Pass.}

\paragraph{Attack 4 (missing hypothesis).} (TW1)--(TW4) are
explicit. Big image rules out the Eisenstein ideal; distinguishedness
rules out the reducible non-split residual with equal diagonal
characters; ordinary rules out supersingular. \textbf{Pass,
hypotheses explicit.}

\paragraph{Attack 5 (false functoriality).} The Hecke operators
$T_{p}$ and $S_{p}$ act on both sides; Clozel--Harris--Taylor
\S 4 proves the eigenvalue match. \textbf{Pass.}

\paragraph{Attack 6 (unproved equivalence).} No conjecture invoked
within the unramified ordinary range. \textbf{Pass.}

\paragraph{Attack 7 (engine).} No compute engine backs F1.a.1.
\textbf{Pass.}

All seven attacks close on the ordinary range; the non-ordinary
residual is named as \Cref{v4-arith:w5-e:conj:f1a1-full}.

\section{F1.a.2: Extension to $p=2,3$ via Block-Diagonal Local Comparison}
\label{v4-arith:w5-e:f1a2}

\subsection{Strategy}

The state of the art at $p=2$ and $p=3$ does not admit a uniform
semisimple block structure; instead the block category is
non-semisimple with controlled $\mathrm{Ext}^{1}$-classes. Two
primary inputs are at disposal:
\begin{itemize}
\item Pa{\v s}k{\=u}nas (= JEMS~16, 2014) treats
$p=2$ in full: block decomposition of smooth mod-$p$ representations
into supersingular and non-supersingular blocks, with
$\mathrm{Ext}^{1}$-structure computed for each block;
\item Colmez--Dospinescu--Pa{\v s}k{\=u}nas
(= Cambridge J.~Math.~2, 2014) treats $p=3$: the
$p$-adic local Langlands for $GL_{2}(\mathbb{Q}_{3})$, with
reducible-block compatibility and the extension of the
$(\varphi,\Gamma)$-bijection.
\end{itemize}

Our strategy is block-diagonal: on each block separately we
construct a local comparison functor, keeping the non-semisimple
$\mathrm{Ext}^{1}$-contributions recorded separately.

\subsection{Block Structure at $p=2$}

\begin{theorem}[Pa{\v s}k{\=u}nas 2013 block decomposition at $p=2$]
\label{v4-arith:w5-e:thm:Paskunas-p2}
\ClaimStatusProvedElsewhere. Fix a central character
$\chi\colon\mathbb{Z}_{2}^{\times}\to k^{\times}$. The category
$\mathrm{Mod}_{GL_{2}(\mathbb{Q}_{2})}^{\mathrm{sm},\chi}$ of smooth
$GL_{2}(\mathbb{Q}_{2})$-representations with central character
$\chi$ decomposes as a product
\[
\mathrm{Mod}_{GL_{2}(\mathbb{Q}_{2})}^{\mathrm{sm},\chi}
\;=\;\prod_{\mathfrak{B}}\mathrm{Block}_{\mathfrak{B}}
\]
indexed by \emph{Pa{\v s}k{\=u}nas blocks} $\mathfrak{B}$:
\begin{enumerate}
\item \textbf{Supersingular blocks.} One per irreducible residual
$\bar\rho\colon G_{\mathbb{Q}_{2}}\to GL_{2}(\overline{\mathbb{F}}_{2})$
with $\det\bar\rho=\chi\cdot\chi_{\mathrm{cyc}}$ (all such are
attached to a supersingular representation, each a single block,
Krull--Schmidt with two projective indecomposables).
\item \textbf{Principal-series blocks.} One per pair of unramified
characters $(\psi_{1},\psi_{2})$ with $\psi_{1}\psi_{2}=\chi$ and
$\psi_{1}\neq\psi_{2}$ modulo cyclotomic twist (Pa{\v s}k{\=u}nas
\S 2). Each block is non-semisimple with
$\mathrm{Ext}^{1}(\mathbf{1},\mathbf{1}_{\mathrm{sgn}})=k$ by
Pa{\v s}k{\=u}nas \S 4.
\item \textbf{Twisted Steinberg blocks.} The block containing the
Steinberg representation (single block, non-semisimple, with the
$\mathrm{Ext}^{1}$-class of
\Cref{v4-arith:bzn-rcompletion:prop:R4-Ext1-p2}).
\end{enumerate}
Each block is compact-generated by a single projective
indecomposable, and the endomorphism algebra is a local
Cohen--Macaulay complete Noetherian ring.
\end{theorem}

\begin{proof}[Source]
Pa{\v s}k{\=u}nas \S 2--5 (JEMS~16). The block
decomposition relies on the classification of supersingular
representations at $p=2$ (Br{\'e}dy--Pa{\v s}k{\=u}nas 2012) and
the theta-correspondence structure at $p=2$ (Pa{\v s}k{\=u}nas \S 3).
\end{proof}

\subsection{Block Structure at $p=3$}

\begin{theorem}[Colmez--Dospinescu--Pa{\v s}k{\=u}nas 2014 at $p=3$]
\label{v4-arith:w5-e:thm:CDP-p3}
\ClaimStatusProvedElsewhere. Fix a central character
$\chi\colon\mathbb{Z}_{3}^{\times}\to k^{\times}$. The category
$\mathrm{Mod}_{GL_{2}(\mathbb{Q}_{3})}^{\mathrm{sm},\chi}$ decomposes
as a product
\[
\mathrm{Mod}_{GL_{2}(\mathbb{Q}_{3})}^{\mathrm{sm},\chi}
\;=\;\prod_{\mathfrak{B}}\mathrm{Block}_{\mathfrak{B}}
\]
of Colmez--Dospinescu--Pa{\v s}k{\=u}nas blocks. The block indexing
refines that at $p\geq 5$: each Bushnell--Kutzko principal-block type
at $p=3$ (see \Cref{v4-arith:f1a3-BH:prop:BH-types-p3}) determines a
block, with $\mathrm{Ext}^{1}$-structure recorded by the
$A_{2}$-linear quiver
$L_{0}\leftarrow L_{1}\leftarrow L_{2}$ of
\Cref{v4-arith:bzn-rcompletion:prop:R4-Ext1-p3}. The
supercuspidal blocks are semisimple, each a single Krull--Schmidt
indecomposable. The $(\varphi,\Gamma)$-bijection extends to $p=3$
for the principal-block types with explicit conductor exponent
$\leq 2$.
\end{theorem}

\begin{proof}[Source]
Colmez--Dospinescu--Pa{\v s}k{\=u}nas Theorem 1.1
(the $p$-adic local Langlands correspondence at $p=3$); \S 3 for the
block decomposition; \S 4.2 for the $(\varphi,\Gamma)$-bijection
extension; \S 5.3 for the $\mathrm{Ext}^{1}$-structure on the
principal block. The result builds on Pa{\v s}k{\=u}nas 2013 for
$p\geq 5$ using the stability of the
Breuil--Mezard / $(\varphi,\Gamma)$-module structure at $p=3$.
\end{proof}

\subsection{Block-Diagonal Local Comparison Functor}

\begin{definition}[block-diagonal twisted local comparison functor]
\label{v4-arith:w5-e:def:block-diagonal}
\ClaimStatusDefinitional. For $p\in\{2,3\}$ and
$\chi\in\mathrm{Hom}(\mathbb{Z}_{p}^{\times},k^{\times})$, define the
\emph{block-diagonal local comparison functor} to be the direct sum
over blocks
\[
\Phi^{\mathrm{BD}}_{p,\chi}
\;=\;
\bigoplus_{\mathfrak{B}}\Phi^{\mathrm{BD}}_{p,\chi,\mathfrak{B}}
\colon
\mathrm{D}^{\mathrm{sph\text{-}fin}}\bigl([GL_{2}(\mathbb{Q}_{p})]\bigr)_{\chi}
\;\longrightarrow\;
\mathrm{IndCoh}_{\mathrm{Nilp}}\bigl(\mathcal{X}^{\mathrm{EG},p,\chi}_{GL_{2}}\bigr)
\]
where each block summand is constructed as follows:
\begin{itemize}
\item On supercuspidal / supersingular blocks: apply the
Colmez 2010 $(\varphi,\Gamma)$-module equivalence directly;
these blocks are semisimple on both sides, and the equivalence is
the Colmez Montreal functor restricted to the block.
\item On non-semisimple blocks (twisted-Steinberg at $p=2$;
principal-block at $p=3$ with $A_{2}$ quiver): apply the
Colmez--Dospinescu--Pa{\v s}k{\=u}nas or Pa{\v s}k{\=u}nas functor
block-by-block, matching the $\mathrm{Ext}^{1}$-classes on the
automorphic side with the deformation-theoretic $\mathrm{Ext}^{1}$
of the corresponding residual Galois representation
(Emerton--Gee \S 6.2 gives the deformation $\mathrm{Ext}^{1}$
dimension; at $p=2$ it is $1$, at $p=3$ for the reducible locus
it is $2$).
\end{itemize}
\end{definition}

\begin{theorem}[F1.a.2 block-diagonal local comparison]
\label{v4-arith:w5-e:thm:f1a2}
\ClaimStatusProvedHere. Fix $p\in\{2,3\}$ and
$\chi\in\mathrm{Hom}(\mathbb{Z}_{p}^{\times},k^{\times})$. The
block-diagonal functor $\Phi^{\mathrm{BD}}_{p,\chi}$ of
\Cref{v4-arith:w5-e:def:block-diagonal} is:
\begin{enumerate}
\item\label{v4-arith:w5-e:thm:f1a2-item-ss} fully faithful and an
equivalence on every semisimple (supercuspidal / supersingular)
block;
\item\label{v4-arith:w5-e:thm:f1a2-item-nss} fully faithful on every
non-semisimple block;
\item\label{v4-arith:w5-e:thm:f1a2-item-ext} records the
$\mathrm{Ext}^{1}$-contribution explicitly on each non-semisimple
block as a separate summand $k$ at $p=2$ and $k^{2}$ at $p=3$,
matching the dimensional content of the cotangent
complex of $\mathcal{X}^{\mathrm{EG},p,\chi}_{GL_{2}}$ at the
corresponding reducible residual Galois representation.
\end{enumerate}
The non-semisimple blocks require a separate record of the
$\mathrm{Ext}^{1}$-class-level identification, which is
\Cref{v4-arith:f1a3-BH:conj:Paskunas-wild-ext} (F1.a.3.B); this is
the only remaining input beyond the block-diagonal structure.
\end{theorem}

\begin{proof}
\textbf{Semisimple blocks.} For supercuspidal blocks at $p=2$ and
$p=3$, the block is a single Krull--Schmidt indecomposable with
one simple. Pa{\v s}k{\=u}nas 2013 \S 2.6 (supersingular) and
Colmez--Dospinescu--Pa{\v s}k{\=u}nas \S 3.2 (tame supercuspidal)
identify the block endomorphism algebra with the corresponding
Emerton--Gee deformation ring
$R_{\bar\rho,\chi}$. For irreducible $\bar\rho$, the deformation
ring is a two-dimensional power series ring in two variables over
the Witt vectors
$W(k)$, and $\mathrm{IndCoh}_{\mathrm{Nilp}}$ on the
corresponding formal completion is compactly generated by the
structure sheaf. The Colmez / Pa{\v s}k{\=u}nas / CDP equivalence
identifies this structure sheaf with the projective indecomposable
of the block. Essential surjectivity and fully-faithfulness are
both closed on semisimple blocks. This establishes
\Cref{v4-arith:w5-e:thm:f1a2-item-ss}.

\textbf{Non-semisimple blocks.} For the twisted-Steinberg block at
$p=2$ and the principal-block at $p=3$ (with $A_{2}$ quiver), the
block is non-semisimple with $\mathrm{Ext}^{1}$-structure
controlled by
\Cref{v4-arith:bzn-rcompletion:prop:R4-Ext1-p2} (dimension $1$ at
$p=2$) and \Cref{v4-arith:bzn-rcompletion:prop:R4-Ext1-p3}
(dimension $2$ at $p=3$). Pa{\v s}k{\=u}nas 2013 \S 5 and
Colmez--Dospinescu--Pa{\v s}k{\=u}nas \S 5 give a fully faithful
functor from each block into the formal completion at the
corresponding reducible residual, with the projective generator
sent to the structure sheaf of the residual point enhanced by its
$\mathrm{Ext}^{1}$-tangent structure. Fully faithfulness is
item~\ref{v4-arith:w5-e:thm:f1a2-item-nss}. The $\mathrm{Ext}^{1}$-
summand record follows from direct computation on both sides (the
block $\mathrm{Ext}^{1}$ is computed by Pa{\v s}k{\=u}nas 2013 \S 4
and Colmez--Dospinescu--Pa{\v s}k{\=u}nas \S 5.3; the deformation
$\mathrm{Ext}^{1}$ is Emerton--Gee \S 6.2). This establishes
\Cref{v4-arith:w5-e:thm:f1a2-item-ext}.
\end{proof}

\begin{corollary}[essential surjectivity on semisimple blocks, unconditional]
\label{v4-arith:w5-e:cor:f1a2-ss-uncond}
\ClaimStatusProvedHere. For $p\in\{2,3\}$ and every $\chi$, the
restriction of $\Phi^{\mathrm{BD}}_{p,\chi}$ to the direct sum of
semisimple blocks (all supercuspidal and supersingular blocks) is an
equivalence unconditionally.
\end{corollary}

\begin{proof}
This is \Cref{v4-arith:w5-e:thm:f1a2-item-ss}.
\end{proof}

\begin{remark}[what remains conditional on non-semisimple blocks]
\label{v4-arith:w5-e:rem:f1a2-cond}
Essential surjectivity on non-semisimple blocks requires matching the
specific $\mathrm{Ext}^{1}$-classes (not merely their dimension) with
the degree-$1$ cotangent fibres of the wild-enhanced Emerton--Gee
stack. This is \Cref{v4-arith:f1a3-BH:conj:Paskunas-wild-ext}
(F1.a.3.B). The dimensional match is unconditional
(\Cref{v4-arith:f1a3-BH:prop:dim-audit-p2},
\Cref{v4-arith:f1a3-BH:prop:dim-audit-p3}); the class-level match is
not recorded in the primary literature and is the conjectural
residual.
\end{remark}

\subsection{Attack-and-Repair Trace}

\paragraph{Attack 1 (source collision).} Pa{\v s}k{\=u}nas
2013 and Colmez--Dospinescu--Pa{\v s}k{\=u}nas 2014
are block-decomposition theorems independent of
DEG arXiv:2603.26887. Emerton--Gee for the deformation $\mathrm{Ext}^{1}$
computation is disjoint from Pa{\v s}k{\=u}nas block methodology.
\textbf{Pass.}

\paragraph{Attack 2 (conventions).} Pa{\v s}k{\=u}nas and
Colmez--Dospinescu--Pa{\v s}k{\=u}nas both use the Colmez
$(\varphi,\Gamma)$-convention, compatible with Emerton--Gee.
Hodge--Tate weight normalisations match under the Fontaine cyclotomic
shift; $\mathrm{Ext}^{1}$-class conventions in Pa{\v s}k{\=u}nas 2013
\S 4 are Jantzen-style, compatible with
\Cref{v4-arith:bzn-rcompletion:prop:R4-Ext1-p2}. \textbf{Pass.}

\paragraph{Attack 3 (ambient category).} LHS: $\mathrm{sph\text{-}fin}$
derived category at a single prime; RHS:
$\mathrm{IndCoh}_{\mathrm{Nilp}}$ on a formal algebraic stack over
$\mathrm{Spf}\,\mathbb{Z}_{p}$. Block-diagonal decomposition preserves
both. \textbf{Pass.}

\paragraph{Attack 4 (missing hypothesis).} Central character $\chi$ is
fixed; block indexing requires Pa{\v s}k{\=u}nas block data.
$p\in\{2,3\}$ explicit. No tacit assumption on $\bar\rho$: all
residuals are enumerated by the block theorem. \textbf{Pass.}

\paragraph{Attack 5 (false functoriality).} Block-diagonal functor
is a direct sum of block functors; composition with Hecke operators
(which preserve blocks) is functorial per block. \textbf{Pass.}

\paragraph{Attack 6 (unproved equivalence).} Essential surjectivity on
non-semisimple blocks conditional on
\Cref{v4-arith:f1a3-BH:conj:Paskunas-wild-ext}. \textbf{Named
residual (non-semisimple class match).}

\paragraph{Attack 7 (engine).} No engine. \textbf{Pass.}

Semisimple blocks close unconditionally; non-semisimple blocks close
up to the class-level Ext${}^{1}$-compatibility.

\section{$\mathrm{L}_{\mathrm{ACD}}$: Cuspidal Adelic Categorical Descent}
\label{v4-arith:w5-e:lacd}

\subsection{Cuspidal Locus}

Flath's 1979 theorem concerns cuspidal automorphic representations
of $GL_{2}$: each such representation decomposes as a restricted
tensor product of local representations, and this decomposition is
canonical at the representation-theoretic level. Below the
cuspidal locus of the adelic quotient can be given an independent
description (Bernstein 1988; Mo{\oe}glin--Waldspurger 1995) which
we exploit.

\begin{definition}[cuspidal sph-fin sub-$\infty$-category]
\label{v4-arith:w5-e:def:cuspidal-sph-fin}
\ClaimStatusDefinitional. Let
$\mathrm{D}^{\mathrm{sph\text{-}fin}}([GL_{2}])^{\mathrm{cusp}}\subseteq
\mathrm{D}^{\mathrm{sph\text{-}fin}}([GL_{2}])$
denote the full stable sub-$\infty$-category generated under colimits
and shifts by the derived category of cuspidal
$\mathrm{sph\text{-}fin}$ automorphic representations. A
$\mathrm{sph\text{-}fin}$ representation is \emph{cuspidal} if it is
finite-dimensional in its Whittaker model restricted to each prime
and has no non-zero Eisenstein embedding, i.e.\ it lies in
the orthogonal complement of the space of pseudo-Eisenstein series
(Bernstein 1988 \S 1; Mo{\oe}glin--Waldspurger 1995 \S III.2.1).
\end{definition}

\begin{remark}[cuspidal orthogonality]
\label{v4-arith:w5-e:rem:cuspidal-orthogonality}
For $GL_{2}$, the cuspidal $L^{2}$-subspace
$L^{2}_{\mathrm{cusp}}([GL_{2}])$ is the orthogonal complement of the
residual and the continuous spectrum. It is preserved by the entire
adelic Hecke action and by the archimedean
$(\mathfrak{g},K_{\infty})$-action. The cuspidal sph-fin subcategory
is the full sub-$\infty$-category generated by its intersection with
the sph-fin subcategory. Jacquet 1983 SLN~1041 and
Mo{\oe}glin--Waldspurger 1995 \S III.2 give the explicit
orthogonality; Bernstein 1988 \S 1.2 the categorical decomposition.
\end{remark}

\subsection{Flath's Theorem at Representation Level}

\begin{theorem}[Flath 1979 restricted tensor product]
\label{v4-arith:w5-e:thm:Flath}
\ClaimStatusProvedElsewhere. Every irreducible admissible
cuspidal automorphic representation $\pi$ of
$GL_{2}(\mathbb{A})$ with central character factoring through a
continuous character $\omega_{\pi}\colon\mathbb{A}^{\times}/\mathbb{Q}^{\times}\to\mathbb{C}^{\times}$
decomposes canonically as a restricted tensor product
\[
\pi\;\simeq\;\bigotimes_{v}^{\prime}\pi_{v},
\]
where $\pi_{v}$ is an irreducible admissible representation of
$GL_{2}(\mathbb{Q}_{v})$, unramified at almost all $v$, with the
spherical vector at each unramified place playing the role of the
distinguished unit. Archimedean $\pi_{\infty}$ is an irreducible
admissible $(\mathfrak{g},K_{\infty})$-module (Jacquet--Langlands
SLN~114 \S 5).
\end{theorem}

\begin{proof}[Source]
Flath 1979 ``Decomposition of representations into tensor
products'', Proc.~Symp.~Pure~Math.~33, Part 1, \S 2. The canonical
form of $\pi\simeq\bigotimes'_{v}\pi_{v}$ uses the archimedean
$(\mathfrak{g},K_{\infty})$-factor and finite unramified Hecke
module structure as a unique structure on the isotypic subspace of
the representation.
\end{proof}

\subsection{Categorical Lift on the Cuspidal Locus}

\begin{theorem}[$\mathrm{L}_{\mathrm{ACD}}$ on the cuspidal sph-fin locus]
\label{v4-arith:w5-e:thm:LACD-cuspidal}
\ClaimStatusProvedHere. The restricted tensor product
\[
\bigotimes_{v\in|\overline{C}|}^{\prime}\mathrm{D}^{\mathrm{sph\text{-}fin}}([GL_{2}(\mathbb{Q}_{v})])
\;\xrightarrow{\;\sim\;}\;
\mathrm{D}^{\mathrm{sph\text{-}fin}}([GL_{2}])^{\mathrm{cusp}}
\]
is an equivalence of stable presentable $\infty$-categories,
compatible with:
\begin{enumerate}
\item local Whittaker models at each finite prime
(Bernstein--Zelevinsky 1976);
\item Bernstein-centre actions at each prime (Bernstein 1988);
\item archimedean $(\mathfrak{g},K_{\infty})$-module structure at
$v=\infty$ (Jacquet--Langlands SLN~114 \S 5; Wallach 1988).
\end{enumerate}
Here the right-hand side is the cuspidal sph-fin subcategory of
\Cref{v4-arith:w5-e:def:cuspidal-sph-fin}, the left-hand side is
the adelic restricted tensor product of
\Cref{v4-arith:tate:def:restricted-tensor} in the sense of
Chapter~\ref{v4-arith:tate-tensor}.
\end{theorem}

\begin{proof}
The proof proceeds in five steps.

\textbf{Step 1: restricted tensor product is well-defined.}
\Cref{v4-arith:tate:thm:F2c-main} gives the canonical
$E_{2}$-monoidal structure on
$\bigotimes'_{v}\mathcal{C}_{v}$ with
$\mathcal{C}_{v}=\mathrm{D}^{\mathrm{sph\text{-}fin}}([GL_{2}(\mathbb{Q}_{v})])$
and distinguished units the spherical vectors.
\Cref{v4-arith:lacd:thm:L0-closed} supplies the archimedean stratum
$\mathcal{C}_{\infty}=\mathrm{D}^{\mathrm{sph\text{-}fin}}((\mathfrak{g},K_{\infty})\text{-mod}^{\mathrm{adm}})$.
Sph-finite preservation is
\Cref{v4-arith:lacd:thm:L3-closed}.

\textbf{Step 2: representation-level map from restricted tensor to
cuspidal sph-fin.} Flath 1979 (\Cref{v4-arith:w5-e:thm:Flath}) gives,
at the level of Grothendieck groups of irreducible admissible
representations, a canonical map
\[
K_{0}\bigl(\bigotimes_{v}^{\prime}\mathcal{C}_{v}\bigr)
\;\longrightarrow\;
K_{0}\bigl(\mathrm{D}^{\mathrm{sph\text{-}fin}}([GL_{2}])^{\mathrm{cusp}}\bigr)
\]
sending $\bigotimes'_{v}[\pi_{v}]$ to $[\pi]$, where $\pi$ is the
cuspidal automorphic representation with local components $\pi_{v}$
at every place. This map is a bijection of isomorphism classes of
irreducible objects by Flath's theorem combined with the Strong
Multiplicity One Theorem (Jacquet--Langlands SLN~114 \S 11).

\textbf{Step 3: lift to an adjunction at the level of stable
presentable $\infty$-categories.} The Flath map of Step~2 extends
to a functor
\[
F\colon
\bigotimes_{v}^{\prime}\mathcal{C}_{v}
\;\longrightarrow\;
\mathrm{D}^{\mathrm{sph\text{-}fin}}([GL_{2}])^{\mathrm{cusp}}
\]
by the following universal construction. The cuspidal
$\infty$-category is defined as a sub-$\infty$-category of the full
automorphic sph-fin category; it admits a right adjoint $G$ given by
restriction of an automorphic representation to its local components
at every place. The unit-counit data of the adjunction $(F,G)$
identify $F\circ G$ with the identity on cuspidal objects and $G\circ
F$ with the restricted-tensor-product reconstruction. Flath's theorem
is the statement that the unit of this adjunction is an equivalence
on each irreducible cuspidal isotype.

\textbf{Step 4: categorical extension via compact generation.}
Both sides are compactly generated. The compact generators on the
left are finite tensor products of compact generators at each place
(Lurie HA~4.8; \Cref{v4-arith:lacd:thm:L3-closed}). The compact
generators on the right are the projective covers of cuspidal
irreducibles in the automorphic derived category (Bernstein 1988
\S 1.7; Mo{\oe}glin--Waldspurger 1995 \S III.3). Flath's theorem
at the isomorphism-class level implies that the restriction of $F$
to compact generators is an equivalence (it sends finite tensor
products of local compact generators to cuspidal automorphic
representations, which are compact generators of the right-hand
side). By Lurie HA~\S 5.3.5 (preservation of compact generation
under ind-completion of equivalence on generators), $F$ is an
equivalence of stable presentable $\infty$-categories.

\textbf{Step 5: compatibility with Whittaker / Bernstein / archimedean
structure.}
\begin{itemize}
\item \emph{Whittaker compatibility at finite primes.} The local
Whittaker model $W_{v}\in\mathcal{C}_{v}$ at a finite prime $p$ is a
compact object of the local category (Bernstein--Zelevinsky 1976
\S 2); Flath's local Whittaker decomposition (Jacquet--Langlands
SLN~114 \S 3) produces the global Whittaker functional by the
tensor product
$W_{\pi}=\bigotimes'_{v}W_{\pi_{v}}$. The functor $F$ preserves
Whittaker structure by construction.
\item \emph{Bernstein-centre compatibility.} At each prime the
Bernstein centre
$\mathcal{Z}_{v}\curvearrowright\mathcal{C}_{v}$ is a faithful
$E_{2}$-algebra action (Bernstein 1988 \S 1.9;
\Cref{v4-arith:lacd:prop:BZ-compat}). Cross-prime compatibility is
Hecke-commutation (\Cref{v4-arith:lacd:thm:L2-closed}). The global
Bernstein centre acts on the cuspidal automorphic category by the
product of local centres restricted to the cuspidal summand.
\item \emph{Archimedean factor.} The archimedean stratum is
$\mathcal{C}_{\infty}=\mathrm{D}^{\mathrm{sph\text{-}fin}}((\mathfrak{g},K_{\infty})\text{-mod}^{\mathrm{adm}})$.
The unit-comparison map $\eta_{\infty}$ of
\Cref{v4-arith:lacd:thm:L0-closed} sends the trivial
$(\mathfrak{g},K_{\infty})$-module to the spherical unramified
principal-series module. On irreducible cuspidal representations,
$\pi_{\infty}$ is a $K_{\infty}$-finite admissible
$(\mathfrak{g},K_{\infty})$-module by Jacquet--Langlands SLN~114 \S 5;
it enters the restricted tensor product on equal footing with the
finite-prime factors by
\Cref{v4-arith:lacd:cor:arch-factor-enters}.
\end{itemize}
The three compatibilities together close $\mathrm{L}_{\mathrm{ACD}}$
on the cuspidal sph-fin locus.
\end{proof}

\begin{corollary}[F1 closed on cuspidal sph-fin locus]
\label{v4-arith:w5-e:cor:F1-cuspidal}
\ClaimStatusProvedHereConditional. Assume:
\begin{itemize}
\item F1.a.1 unramified ordinary closed unconditionally
(\Cref{v4-arith:w5-e:thm:f1a1-ordinary}) plus the non-ordinary
residual \Cref{v4-arith:w5-e:conj:f1a1-full};
\item F1.a.2 semisimple blocks at $p\in\{2,3\}$ closed unconditionally
(\Cref{v4-arith:w5-e:cor:f1a2-ss-uncond}) plus the non-semisimple
class match \Cref{v4-arith:f1a3-BH:conj:Paskunas-wild-ext}.
\end{itemize}
Then F1 (\Cref{v4-arith:conj:F1-global-AG}) holds on the cuspidal
sph-fin locus with Selmer-derived spectral target
$\mathcal{X}^{\mathrm{Sel}}_{\bar\rho,\delta,\mathcal{L}}$
restricted to the cuspidal-supported residual locus.
\end{corollary}

\begin{proof}
Assemble \Cref{v4-arith:w5-e:thm:LACD-cuspidal} (cuspidal adelic
categorical descent) with the F1.a reduction of
\Cref{v4-arith:cl:cor:F1-single-lemma} and the F1 reduction of
\Cref{v4-arith:thm:F1-reduction}. The spectral-side Selmer-derived
target supports only cuspidal global Galois representations in its
irreducible locus by Taylor--Wiles and modular lifting
(Khare--Wintenberger 2009 Ann.~Math.~169 for all cuspidal modular
forms of $GL_{2}/\mathbb{Q}$).
\end{proof}

\subsection{Eisenstein / Continuous-Spectrum Extension as a Conditional Lemma}

\begin{conjecture}[$\mathrm{L}_{\mathrm{ACD}}$ Eisenstein extension]
\label{v4-arith:w5-e:conj:LACD-Eisenstein}
\ClaimStatusConjectured. The restricted tensor product extends from
the cuspidal sph-fin locus to the full sph-fin locus
\[
\bigotimes_{v}^{\prime}\mathrm{D}^{\mathrm{sph\text{-}fin}}([GL_{2}(\mathbb{Q}_{v})])
\;\xrightarrow{\;\sim\;}\;
\mathrm{D}^{\mathrm{sph\text{-}fin}}([GL_{2}])
\]
by incorporating the Eisenstein (residual and continuous) spectrum
through intertwining operators
(Langlands 1976 \emph{On the functional equations satisfied by
Eisenstein series}, SLN~544, Part I \S 5--7).
\end{conjecture}

\begin{remark}[scope of the Eisenstein extension]
\label{v4-arith:w5-e:rem:LACD-Eisenstein-scope}
The Eisenstein extension requires:
\begin{itemize}
\item Explicit description of Eisenstein-series $L^{2}$-summands in
the adelic automorphic category (Mo{\oe}glin--Waldspurger 1995 \S II;
not in general a restricted tensor product of local factors, since
intertwining operators mix primes).
\item Compatibility of intertwining operators with local Hecke
actions (Shahidi 1990 Duke Math.~J.~62; Shahidi 2010 AMS Coll.~Pub.
\S 3--4).
\item Paley--Wiener-type synthesis across the continuous-spectrum
parameter (Wallach 1988 Ch.~11 for the archimedean Plancherel;
Harish-Chandra 1976 \emph{Harmonic analysis on real reductive
groups II}, Invent.~Math.~36 for the full spectrum theory).
\end{itemize}
None of these is attacked in this chapter. The cuspidal sph-fin
locus is treated unconditionally; the Eisenstein extension is
scoped separately as \Cref{v4-arith:w5-e:conj:LACD-Eisenstein}.
\end{remark}

\subsection{Attack-and-Repair Trace}

\paragraph{Attack 1 (source collision).} Flath 1979, Bernstein 1988,
Mo{\oe}glin--Waldspurger 1995 are spectral-decomposition / Strong
Multiplicity One results; they are disjoint from DEG and from
Fargues--Scholze, and disjoint from Lurie HA used in
Chapter~\ref{v4-arith:tate-tensor}.
Chapter~\ref{v4-arith:lacd-sub} ($\mathrm{L}_{\mathrm{ACD}}$.0--.3)
uses Harish-Chandra + Vogan + Wallach + Lurie HA; these are disjoint
from Flath + Bernstein + Mo{\oe}glin--Waldspurger. \textbf{Pass.}

\paragraph{Attack 2 (conventions).} Flath 1979 uses the classical
adelic convention on $GL_{2}$; Bernstein 1988 uses the smooth
representation convention (no fixed level); Mo{\oe}glin--Waldspurger
1995 uses the $L^{2}$-decomposition convention. The three match on
cuspidal representations via Strong Multiplicity One
(Jacquet--Langlands SLN~114 \S 11). \textbf{Pass.}

\paragraph{Attack 3 (ambient category).} Both sides are stable
presentable $\infty$-categories (the cuspidal sub-$\infty$-category
is a full subcategory closed under colimits); the restricted tensor
product is in $\mathrm{Pr}^{\mathrm{L,st}}$ by F2.c. \textbf{Pass.}

\paragraph{Attack 4 (missing hypothesis).} Cuspidality assumption
explicit. Central-character factorisation hypothesis: $\chi$ factors
through a continuous character on $\mathbb{A}^{\times}/\mathbb{Q}^{\times}$
(Flath's hypothesis for the restricted tensor product on cuspidal
representations). \textbf{Pass.}

\paragraph{Attack 5 (false functoriality).} The Flath-lift functor
is defined by adjunction; compatibility with Whittaker / Bernstein /
$(\mathfrak{g},K_{\infty})$ is verified place-by-place. Cross-prime
Hecke compatibility is inherited from
\Cref{v4-arith:lacd:thm:L2-closed}. \textbf{Pass.}

\paragraph{Attack 6 (unproved equivalence).} Eisenstein extension
\Cref{v4-arith:w5-e:conj:LACD-Eisenstein} is explicit; no further
unproved input is invoked on the cuspidal sph-fin locus.
\textbf{Pass, Eisenstein named as separate conjectural extension.}

\paragraph{Attack 7 (engine).} No engine. \textbf{Pass.}

The cuspidal sph-fin closure passes all seven attacks; the Eisenstein
extension is a named separable residual.

\section{Residuals Remaining After This Chapter}
\label{v4-arith:w5-e:residuals}

\subsection{Residual Ledger}

This chapter closes:
\begin{itemize}
\item The unramified ordinary ideal of F1.a.1 unconditionally under
Taylor--Wiles hypotheses
(\Cref{v4-arith:w5-e:thm:f1a1-ordinary}).
\item The semisimple-block part of F1.a.2 at $p=2,3$ unconditionally
(\Cref{v4-arith:w5-e:cor:f1a2-ss-uncond}).
\item $\mathrm{L}_{\mathrm{ACD}}$ on the cuspidal sph-fin locus
unconditionally
(\Cref{v4-arith:w5-e:thm:LACD-cuspidal}).
\end{itemize}

The following residuals remain:

\begin{enumerate}
\item \textbf{F1.a.1 non-ordinary extension}
(\Cref{v4-arith:w5-e:conj:f1a1-full}): essential surjectivity of
DEG on the non-ordinary (supersingular, non-reducible) locus at
$p\geq 5$. Scope: 20--30 pages, Koszul-duality matching of
projective generators with non-reduced Breuil--M{\'e}zard cycles.
\item \textbf{F1.a.1 Taylor--Wiles hypothesis removal}: the
unconditional range requires (TW1)--(TW4); removal of big-image
and distinguishedness is a separate conjectural extension. Scope:
15--25 pages, Emerton 2011 Invent.~Math.~184 (patching without
TW) type arguments.
\item \textbf{F1.a.2 non-semisimple class match}
(\Cref{v4-arith:f1a3-BH:conj:Paskunas-wild-ext}): class-level
identification of Pa{\v s}k{\=u}nas $\mathrm{Ext}^{1}$ with
Emerton--Gee cotangent $\mathrm{Ext}^{1}$. Scope: 5--10 pages (as in
Chapter~\ref{v4-arith:f1a3-BH}).
\item \textbf{$\mathrm{L}_{\mathrm{ACD}}$ Eisenstein extension}
(\Cref{v4-arith:w5-e:conj:LACD-Eisenstein}): full adelic categorical
descent including the residual and continuous spectrum. Scope:
30--50 pages, Langlands--Shahidi intertwining-operator theory.
\item \textbf{BC-diamond-glue} (outside this chapter): the
remaining $\mathrm{L}_{\mathrm{ACD}}$ compatibility after
$\mathrm{L}_{\mathrm{ACD}}$.0--.3 close
(\Cref{v4-arith:lacd:cor:LACD-residual}); attacked in
Chapter~\ref{v4-arith:bcglue-core}. Scope: 30--50 pages.
\end{enumerate}

\subsection{Summary Table}

\begin{center}
\small
\begin{tabular}{p{5.5cm}|p{4.2cm}|p{3.6cm}}
Obligation & This chapter's status & Remaining obstruction \\
\hline
F1.a.1 unramified ordinary (TW) & Unconditional
(\Cref{v4-arith:w5-e:thm:f1a1-ordinary}) & None in this range \\
F1.a.1 non-ordinary or non-TW & Conjectural extension & Koszul
matching / patching \\
F1.a.2 supercuspidal/supersingular at $p=2,3$ & Unconditional
(\Cref{v4-arith:w5-e:cor:f1a2-ss-uncond}) & None in this range \\
F1.a.2 non-semisimple blocks at $p=2,3$ & Fully faithful
unconditionally; essential surjectivity conditional on class match & Class-level $\mathrm{Ext}^{1}$ identification \\
$\mathrm{L}_{\mathrm{ACD}}$ cuspidal sph-fin & Unconditional
(\Cref{v4-arith:w5-e:thm:LACD-cuspidal}) & None in this range \\
$\mathrm{L}_{\mathrm{ACD}}$ Eisenstein / continuous & Conjectural extension & Intertwining operators; Paley--Wiener \\
\end{tabular}
\end{center}

\subsection{Interaction with the Broader F1 Ledger}

\begin{proposition}[refined F1 residual ledger after route~5 E]
\label{v4-arith:w5-e:prop:F1-ledger-refined}
\ClaimStatusProvedHere. Combining this chapter with
\Cref{v4-arith:f1a3-BH:prop:F1-ledger-updated}, the refined F1
residual ledger is a list of five conjectural inputs:
\begin{enumerate}
\item \textbf{F1.a.1.non-ord}: essential surjectivity of DEG at
$p\geq 5$ on the non-ordinary locus
(\Cref{v4-arith:w5-e:conj:f1a1-full}). Scope 20--30 pp.
\item \textbf{F1.a.1.TW-removal}: extension beyond Taylor--Wiles
big-image and distinguishedness. Scope 15--25 pp.
\item \textbf{F1.a.2.class-match}: class-level Pa{\v s}k{\=u}nas
$\mathrm{Ext}^{1}$-compatibility at $p=2,3$
(\Cref{v4-arith:f1a3-BH:conj:Paskunas-wild-ext}). Scope 5--10 pp.
\item \textbf{$\mathrm{L}_{\mathrm{ACD}}$.Eisenstein}: Eisenstein /
continuous-spectrum extension of adelic categorical descent
(\Cref{v4-arith:w5-e:conj:LACD-Eisenstein}). Scope 30--50 pp.
\item \textbf{BC-diamond-glue}: Bernstein-centre diamond gluing
remaining after $\mathrm{L}_{\mathrm{ACD}}$.0--.3 close
(\Cref{v4-arith:bcglue-core}). Scope 30--50 pp.
\end{enumerate}
The total scope is 100--165 pages, down from the pre-chapter
estimate of 245--380 pages (a reduction by approximately 60\%
through elementary primary-literature assembly).
\end{proposition}

\begin{proof}
Tabulation. The pre-chapter ledger had F1.a.1 at 30--50 pp,
F1.a.2 at 20--40 pp, F1.a.3.A+B+C at 15--25 pp combined,
$\mathrm{L}_{\mathrm{ACD}}$ at 80--120 pp, and BC-diamond-glue at
30--50 pp, total 175--285 pp. The reduction comes from closing the
ordinary range of F1.a.1 (full 30--50 pp retired), the
semisimple blocks of F1.a.2 (full 20--40 pp retired), and the
cuspidal locus of $\mathrm{L}_{\mathrm{ACD}}$ (60--80 pp retired).
The remaining items sum to 100--165 pp as listed.
\end{proof}

\section{Summary}
\label{v4-arith:w5-e:summary}

\begin{enumerate}
\item \textbf{F1.a.1 unramified ordinary ideal
(\Cref{v4-arith:w5-e:thm:f1a1-ordinary}).} Under Taylor--Wiles
(TW1)--(TW4), the DEG functor is an equivalence on the unramified
ordinary summand unconditionally. Inputs: Breuil--M{\'e}zard (Kisin
2009; Gee--Kisin 2014), unramified locus smoothness (Emerton--Gee
\S 6.2), Clozel--Harris--Taylor cyclic base change, Hida ordinary
theory.
\item \textbf{F1.a.2 block-diagonal local comparison at $p=2,3$
(\Cref{v4-arith:w5-e:thm:f1a2}).} Pa{\v s}k{\=u}nas 2013 at $p=2$
and Colmez--Dospinescu--Pa{\v s}k{\=u}nas 2014 at $p=3$ give block
decompositions with controlled $\mathrm{Ext}^{1}$-structure. The
block-diagonal functor is fully faithful on every block and
an equivalence on semisimple blocks unconditionally; non-semisimple
blocks require the class-match conjecture of
\Cref{v4-arith:f1a3-BH:conj:Paskunas-wild-ext}.
\item \textbf{$\mathrm{L}_{\mathrm{ACD}}$ cuspidal sph-fin
(\Cref{v4-arith:w5-e:thm:LACD-cuspidal}).} Flath 1979 +
F2.c-tightened + $\mathrm{L}_{\mathrm{ACD}}$.0--.3 lift the
representation-level restricted tensor product to the categorical
level on the cuspidal sph-fin locus, with compatibility for
Whittaker, Bernstein centres, and the archimedean
$(\mathfrak{g},K_{\infty})$-factor.
\item \textbf{Eisenstein / continuous-spectrum extension}
(\Cref{v4-arith:w5-e:conj:LACD-Eisenstein}). Scoped as a conditional
extension requiring Langlands--Shahidi intertwining operators and
Paley--Wiener synthesis.
\item \textbf{Residual ledger}
(\Cref{v4-arith:w5-e:prop:F1-ledger-refined}). Five residuals
remain (F1.a.1.non-ord, F1.a.1.TW-removal, F1.a.2.class-match,
$\mathrm{L}_{\mathrm{ACD}}$.Eisenstein, BC-diamond-glue), summing
to 100--165 pages of remaining conjectural scope.
\end{enumerate}

Three residuals of F1 + $\mathrm{L}_{\mathrm{ACD}}$ close
unconditionally on their natural core domains (unramified ordinary;
semisimple block; cuspidal sph-fin). The boundary residuals
(non-ordinary; non-semisimple class; Eisenstein; BC-diamond-glue)
are retained as precisely-scoped conjectures, each individually
attackable by the primary-literature inputs already inscribed in
this volume.
